\subsection{Arithmetic and undecidability for groups}
\label{sec:p02-group-theory-undecidability}

The decision problem here takes a first-order sentence in the group
signature \(\{\cdot,{}^{-1},1\}\) and asks whether it is true in every
group. A finitely axiomatized subclass suffices to transfer the arithmetic
obstruction to this problem. We use the centre and multiplication formulas
of \cref{def:p03-w03-p03-074-ring-formulas}, with the commutator convention
of \cref{def:p03-w03-p03-074-conventions}.

\begin{hbdefinition}{def:p02-finite-ring-group-condition}{A finite ring condition}
Abbreviate \(\mathrm{Mul}(a,b;z,w,v)\) by \(M_{a,b}(z,w,v)\).
Let \(\Theta_0(a,b)\) be the conjunction of
\[
 \forall u\,\forall v\,\forall w\;[[u,v],w]=1,
\]
\[
 \forall z\bigl(([z,a]=1\land[z,b]=1)
                         \leftrightarrow\mathrm{Ctr}(z)\bigr),
\]
and \(\beta(a,b)\) from \cref{def:p03-w03-p03-074-ring-formulas}.
Define \(\Theta_{\mathrm{ring}}(a,b)\) by adjoining the two formulas
\[
 \forall z\,\forall w\,\forall v\;
   \bigl(M_{a,b}(z,w,v)\to M_{a,b}(w,z,v)\bigr)
\]
and
\[
 \begin{gathered}
 \forall z\,\forall w\,\forall u\,\forall r\,\forall s\,\forall t\\
 \begin{aligned}
 \bigl(&M_{a,b}(z,w,r)\land M_{a,b}(r,u,t)\\
       &\land M_{a,b}(w,u,s)\longrightarrow M_{a,b}(z,s,t)\bigr).
 \end{aligned}
 \end{gathered}
\]
All displayed predicates are the previously defined group formulas;
expanding their abbreviations gives one finite formula in the group
signature, with only \(a,b\) free.
\end{hbdefinition}

\begin{hblemma}{thm:p02-finite-ring-group-condition}{proved=annex:p02.finite-ring-group-condition}{The interpreted ring}
If \(G\models\Theta_{\mathrm{ring}}(a,b)\), the data \(I\) of
\cref{def:p03-w03-p03-074-ring-formulas} define a commutative ring with
identity on \(Z(G)\): zero is the group identity, one is \([a,b]\),
addition is group multiplication, and multiplication has graph
\(M_{a,b}\). Conversely, for every commutative ring \(R\) with identity,
the standard pair \((e_1,e_2)\) in \(\mathrm H(R)\) satisfies
\(\Theta_{\mathrm{ring}}\), and its interpreted ring is isomorphic to
\(R\). The zero ring is allowed at this stage.
\end{hblemma}

\begin{hbdefinition}{def:p02-finite-group-arithmetic}{The auxiliary group theory}
Let \(D(z)\) and \(J(\chi)=\chi^{\mathbb Z}\) denote exactly the
four-square formula and arithmetic-to-ring translation of
\cref{def:p02-w02c-p02-003-ring-coding}. This is a syntactic translation;
its ring formula can be evaluated in rings other than \(\mathbb Z\).
Write \(D(t)\) by capture-avoiding substitution, renaming its bound
variables freshly. Put
\[
 \begin{aligned}
 \mathrm{Cl}_D := {}&D(0)\land D(1)\\
 &\land\forall z\,\forall w\bigl((D(z)\land D(w))\to
                   (D(z+w)\land D(z\cdot w))\bigr).
 \end{aligned}
\]
Let \(A_{\mathsf Q}\) be the seven sentence axioms displayed in
\cref{tab:p05-q-pa-axioms}, and define the single ring sentence
\[
 \vartheta=\mathrm{Cl}_D\land
                         \bigwedge_{\alpha\in A_{\mathsf Q}}J(\alpha).
\]
This conjunction has seven arithmetic-axiom instances, not an induction
scheme. Translate it by the fixed ring interpretation \(I\), and set
\[
 \Theta(a,b)=\Theta_{\mathrm{ring}}(a,b)\land\vartheta^I(a,b).
\]
The theory \(V_G\) consists of the group axioms and
\(\Theta(c_x,c_y)\), in the group signature expanded by the two constants
\(c_x,c_y\) already used in \(I\). Its arithmetic interpretation has
domain \(\mathrm{Ctr}(z)\land D(z)^I\), literal equality, zero graph
\(v=1\), successor graph \(v=u[c_x,c_y]\), addition graph \(w=uv\),
and multiplication graph \(M_{c_x,c_y}(u,v,w)\), on that domain.
Let \(K\) be the recursive graph translation for these data, with
quantifiers restricted to this domain. For every arithmetic sentence
\(\psi\), the proof establishes
\[
 V_G\derives K(\psi)\leftrightarrow
                         \bigl(J(\psi)\bigr)^I(c_x,c_y).
\]
The constants are part of the target signature, so the interpretation
uses no unnamed parameters.
\end{hbdefinition}

\begin{hbtheorem}{thm:p02-finite-group-arithmetic}{proved=annex:p02.finite-group-arithmetic}{The arithmetic interpretation}
The theory \(V_G\) is finitely axiomatized and consistent. The composite
translation \(K\) gives an effective one-dimensional relative
interpretation of \(\mathsf Q\) in \(V_G\), with literal equality.
Its domain consists of the central elements satisfying the interpreted
four-square formula. The group \(\mathrm H(\mathbb Z)\), with
\(c_x=e_1,c_y=e_2\), is a model of \(V_G\), and its interpreted arithmetic is the
standard structure of natural numbers.
\end{hbtheorem}

\begin{hbtheorem}{thm:p02-group-theory-undecidable}{proved=annex:p02.group-theory-undecidable}{Undecidability for groups}
The set of first-order group sentences true in every group is computably
enumerable and undecidable. The auxiliary theory \(V_G\) is essentially
undecidable. Essential undecidability is asserted of this auxiliary theory,
not of the unrestricted group axioms.
\end{hbtheorem}

The conclusion concerns first-order theory validity, not the word problem
for an individual group presentation. The two parameter constants in
\(V_G\) are eliminated by universal quantification in the decision
reduction, so the input sentences for the group-theory problem contain
neither those constants nor additional relation symbols.
